\documentclass[UTF8,12pt]{ctexart}
\usepackage[a4paper,margin=24mm]{geometry}
\usepackage{amsmath,amssymb,bm,graphicx,adjustbox}
\newcommand{\fitmath}[1]{\begin{adjustbox}{max width=\linewidth}$\displaystyle #1$\end{adjustbox}}
\setlength{\parindent}{0pt}
\setlength{\parskip}{.7em}
\sloppy
\allowdisplaybreaks
\begin{document}
\section*{2004 年考研数学三 · 参考答案}
\subsection*{2004年真题参考答案}

\subsection*{一、填空题}

（1）\(\displaystyle 1;-4\)。（2）\(\displaystyle -\dfrac{\displaystyle g^{\prime}(v)}{\displaystyle [g(v)]^2}\)。（3）\(\displaystyle -\dfrac{\displaystyle 1}{\displaystyle 2}\)。（4）\(\displaystyle 2\)。（5）\(\displaystyle \dfrac{\displaystyle 1}{\displaystyle e}\)。（6）\(\displaystyle \sigma^2\)。


\subsection*{二、选择题}

（7）A。（8）D。（9）C。（10）B。（11）D。（12）D。（13）B。（14）C。


\subsection*{三、解答题}

（15）\(\displaystyle \dfrac{\displaystyle 4}{\displaystyle 3}\)。 （16）\(\displaystyle \dfrac{\displaystyle 16}{\displaystyle 9}(3\pi-2)\)。 （17）证明略。


（18）（I）\(\displaystyle E_d=\dfrac{\displaystyle P}{\displaystyle 20-P}\)。（II）当 \(\displaystyle 10<P<20\) 时，价格降低反而使收益增加。


（19）（I）\(\displaystyle y^{\prime}=xy+\dfrac{\displaystyle x^3}{\displaystyle 2}\)，初始条件为 \(\displaystyle y(0)=0\)。（II）\(\displaystyle S(x)=e^{x^2/2}-\dfrac{\displaystyle x^2}{\displaystyle 2}-1\;(-\infty<x<+\infty)\)。


（20）（I）当 \(\displaystyle a=0\)，\(\displaystyle b\) 为任意常数时，\(\displaystyle \beta\) 不能由 \(\displaystyle \alpha_1,\allowbreak \alpha_2,\allowbreak \alpha_3\) 线性表示。（II）当 \(\displaystyle a\ne0\) 且 \(\displaystyle a\ne b\) 时，\(\displaystyle \beta\) 能唯一地线性表示为 \(\displaystyle \beta=(1-\dfrac{\displaystyle 1}{\displaystyle a})\alpha_1+\dfrac{\displaystyle 1}{\displaystyle a}\alpha_2\)。（III）当 \(\displaystyle a=b\ne0\) 时，\(\displaystyle \beta\) 可线性表示但不唯一：\(\displaystyle \beta=(1-\dfrac{\displaystyle 1}{\displaystyle a})\alpha_1+(\dfrac{\displaystyle 1}{\displaystyle a}+k)\alpha_2+k\alpha_3\)，其中 \(\displaystyle k\) 为任意常数。


（21）（I）当 \(\displaystyle b=0\) 或 \(\displaystyle n=1\) 时，\(\displaystyle A=E\)，\(\displaystyle A\) 的特征值 \(\displaystyle \lambda_1=\cdots=\lambda_n=1\)，任意非零列向量均为 \(\displaystyle A\) 的特征向量。当 \(\displaystyle b\ne0\) 且 \(\displaystyle n\ge2\) 时，\(\displaystyle \lambda_1=1+(n-1)b\)，\(\displaystyle \lambda_2=\cdots=\lambda_n=1-b\)；前者的特征向量为 \(\displaystyle k_1(1,\allowbreak 1,\allowbreak \ldots,\allowbreak 1)^{\mathrm T}\)，后者的特征向量为 \(\displaystyle k_2(1,\allowbreak -1,\allowbreak 0,\allowbreak \ldots,\allowbreak 0)^{\mathrm T}+\cdots+k_n(1,\allowbreak 0,\allowbreak \ldots,\allowbreak 0,\allowbreak -1)^{\mathrm T}\)，其中 \(\displaystyle k_i\) 不全为零。


（21）（II）\(\displaystyle P=\begin{pmatrix}1&1&1&\cdots&1\\1&-1&0&\cdots&0\\1&0&-1&\cdots&0\\\vdots&\vdots&\vdots&&\vdots\\1&0&\cdots&0&-1\end{pmatrix}\)。


（22）（I）\(\displaystyle (X,\allowbreak Y)\) 的概率分布为 \(\displaystyle \begin{array}{c|cccc}(X,\allowbreak Y)&(0,\allowbreak 0)&(0,\allowbreak 1)&(1,\allowbreak 0)&(1,\allowbreak 1)\\\hline P&\dfrac{\displaystyle 2}{\displaystyle 3}&\dfrac{\displaystyle 1}{\displaystyle 12}&\dfrac{\displaystyle 1}{\displaystyle 6}&\dfrac{\displaystyle 1}{\displaystyle 12}\end{array}\)。（II）\(\displaystyle \rho_{XY}=\dfrac{\displaystyle 1}{\displaystyle \sqrt{15}}\)。（III）\(\displaystyle P\{Z=0\}=\dfrac{\displaystyle 2}{\displaystyle 3}\)，\(\displaystyle P\{Z=1\}=\dfrac{\displaystyle 1}{\displaystyle 4}\)，\(\displaystyle P\{Z=2\}=\dfrac{\displaystyle 1}{\displaystyle 12}\)。


（23）（I）\(\displaystyle \hat\beta=\dfrac{\displaystyle \bar X}{\displaystyle \bar X-1}\)。（II）\(\displaystyle \hat\beta=\dfrac{\displaystyle n}{\displaystyle \sum_{i=1}^n\ln X_i}\)。（III）\(\displaystyle \hat\alpha=\displaystyle \min\{X_1,\allowbreak X_2,\allowbreak \ldots,\allowbreak X_n\}\)。

\end{document}
